\section{The Circuit Pipeline} \label{sec:circuit-pipeline}

Writing a ZK smart contract on Cardano is not a single-language problem. It involves a chain of tools that each play a different role:

\begin{enumerate}
  \item \textbf{Circom} --- a domain-specific language for describing arithmetic circuits.
  \item \textbf{snarkjs} --- a JavaScript library that runs the trusted setup ceremony, compiles witness calculations, and generates / verifies Groth16 proofs.
  \item \textbf{Aiken} --- the on-chain language where the Cardano validator lives and where the ZK proof is verified against the circuit's public verification key.
\end{enumerate}

The high-level flow is:

\begin{center}
\begin{tabular}{ccccc}
  Circom circuit & $\longrightarrow$ & snarkjs trusted setup & $\longrightarrow$ & verification key \\
  & & $\downarrow$ & & $\downarrow$ \\
  & & Groth16 proof & $\longrightarrow$ & Aiken validator
\end{tabular}
\end{center}

The circuit defines the \emph{rules}. The trusted setup ceremony produces the public parameters that let anyone prove and verify statements about those rules. The proof is generated off-chain by the user. The Aiken validator checks the proof on-chain --- a single, fast pairing check using the BLS12-381 primitives introduced in Plutus V3.

% ─────────────────────────────────────────────────────────────────────────────
\section{Writing a Circuit in Circom} \label{sec:circom}

Circom is a language for expressing computations as \textbf{arithmetic circuits} --- networks of addition and multiplication gates over a finite field. Every statement you want to prove must be reducible to a set of polynomial constraints called an \textbf{R1CS} (Rank-1 Constraint System).

\subsection*{Installing the Tools}

Install snarkjs globally:

\begin{verbatim}
npm install -g snarkjs
\end{verbatim}

Install Circom by following the official guide at \texttt{docs.circom.io/getting-started/installation}. On Linux/macOS this compiles from Rust source:

\begin{verbatim}
git clone https://github.com/iden3/circom.git
cd circom
cargo build --release
cargo install --path circom
\end{verbatim}

Install the circomlib circuit library (contains Poseidon and other primitives):

\begin{verbatim}
npm install -D circomlib
\end{verbatim}

\subsection*{The Janus Wallet Authentication Circuit}

Janus Wallet\footnote{\texttt{https://github.com/leobel/janus-wallet}} is the most complete open-source ZK project on Cardano. It uses a circuit that proves a user knows their password without ever sending it on-chain. The circuit is built on the \textbf{Poseidon} hash function --- a ZK-friendly hash designed to be cheap to prove inside arithmetic circuits (unlike SHA-256, which is extremely expensive to circuit-encode).

The circuit has five \textbf{public} inputs (visible on-chain) and one \textbf{private} input (the secret that never leaves the user's device):

\begin{verbatim}
pragma circom 2.1.6;

include "../../../node_modules/circomlib/circuits/poseidon.circom";

template Authenticate() {
    signal input userId;        // Public: username identifier
    signal input credentialHash; // Public: Poseidon(pwd, userId)
    signal input challenge;      // Public: nonce from the smart contract
    signal input challengeFlag;  // Public: overflow flag for challenge
    signal input circuitHash;    // Public: Poseidon(pwd, userId, credentialHash,
                                 //                  challenge, challengeFlag)
    signal input pwd;            // Private: user password

    // Step 1: verify credentialHash = Poseidon(pwd, userId)
    component p = Poseidon(2);
    p.inputs[0] <== pwd;
    p.inputs[1] <== userId;
    credentialHash === p.out;

    // Step 2: verify circuitHash = Poseidon(pwd, userId, credentialHash,
    //                                        challenge, challengeFlag)
    component poseidon = Poseidon(5);
    poseidon.inputs[0] <== pwd;
    poseidon.inputs[1] <== userId;
    poseidon.inputs[2] <== credentialHash;
    poseidon.inputs[3] <== challenge;
    poseidon.inputs[4] <== challengeFlag;
    circuitHash === poseidon.out;
}

component main {public [userId, credentialHash, challenge,
                         challengeFlag, circuitHash]} = Authenticate();
\end{verbatim}

The \texttt{===} operator asserts a constraint: if the constraint is violated, no valid proof can be produced. The \texttt{<==} operator assigns and constrains in one step.

The \texttt{challenge} and \texttt{challengeFlag} fields are critical for replay protection. The smart contract changes the challenge after every successful spend, so a proof generated for one transaction cannot be reused in another.

\subsection*{Compiling the Circuit}

\begin{verbatim}
circom --r1cs --wasm --c --sym --inspect \
       --prime bls12381 circuit.circom -o .
\end{verbatim}

\begin{remark}
Run this command from the directory where your \texttt{include} path can reach the circomlib node module. For Janus Wallet this is \texttt{src/zkproof/janus-wallet/}.
\end{remark}

The flags produce:

\begin{itemize}
  \item \texttt{circuit.r1cs} --- the constraint system in binary format.
  \item \texttt{circuit\_js/circuit.wasm} --- WebAssembly witness calculator.
  \item \texttt{circuit.sym} --- symbol table for debugging.
\end{itemize}

The \texttt{--prime bls12381} flag targets the BLS12-381 curve, which matches Cardano's on-chain primitives.

Inspect the compiled circuit:

\begin{verbatim}
snarkjs r1cs info circuit.r1cs
\end{verbatim}

For the Janus circuit you should see:

\begin{verbatim}
[INFO] snarkJS: Curve: bls12-381
[INFO] snarkJS: # of Wires: 1357
[INFO] snarkJS: # of Constraints: 1352
[INFO] snarkJS: # of Private Inputs: 1
[INFO] snarkJS: # of Public Inputs: 5
[INFO] snarkJS: # of Outputs: 0
\end{verbatim}

% ─────────────────────────────────────────────────────────────────────────────
\section{The Trusted Setup Ceremony} \label{sec:trusted-setup}

Groth16 requires a two-phase trusted setup. Phase 1 (Powers of Tau) is circuit-independent and can be reused across projects. Phase 2 is circuit-specific and ties the parameters to your exact R1CS.

\subsection*{Phase 1: Powers of Tau}

\paragraph{1. Start a new ceremony.}

\begin{verbatim}
snarkjs powersoftau new bls12381 14 pot14_0000.ptau -v
\end{verbatim}

The second argument (\texttt{14}) is the power of two of the maximum number of constraints: $2^{14} = 16{,}384$. Choose the smallest power that fits your circuit to keep file sizes manageable. The maximum supported is $2^{28} \approx 268$ million constraints.

\paragraph{2--3. Contribute randomness.}

Multiple contributors make the ceremony secure: the parameters are safe as long as \emph{at least one} contributor honestly discarded their secret after contributing.

\begin{verbatim}
snarkjs powersoftau contribute pot14_0000.ptau pot14_0001.ptau \
    --name="First contribution" -v

snarkjs powersoftau contribute pot14_0001.ptau pot14_0002.ptau \
    --name="Second contribution" -v -e="some random text"
\end{verbatim}

The \texttt{-e} flag provides entropy inline, making the step non-interactive --- useful in scripts.

\paragraph{4. (Optional) Third-party contribution.}

You can mix contributions from different software:

\begin{verbatim}
snarkjs powersoftau export challenge pot14_0002.ptau challenge_0003

snarkjs powersoftau challenge contribute bls12381 challenge_0003 \
    response_0003 -e="some random text"

snarkjs powersoftau import response pot14_0002.ptau response_0003 \
    pot14_0003.ptau -n="Third contribution name"
\end{verbatim}

\paragraph{5. Verify the ceremony so far.}

\begin{verbatim}
snarkjs powersoftau verify pot14_0003.ptau
\end{verbatim}

A clean run ends with \texttt{Powers Of tau OK!}

\paragraph{6. Apply a random beacon.}

A random beacon finalises Phase 1. It is a delayed hash function evaluated on a publicly known, high-entropy value --- for example, the hash of a specific Cardano block, or a national lottery result announced after the ceremony began.

\begin{verbatim}
snarkjs powersoftau beacon pot14_0003.ptau pot14_beacon.ptau \
    3c692631872308ff1f9fed102b940aecee160def9c2478b50b977091 \
    10 -n="Final Beacon"
\end{verbatim}

The \texttt{10} argument means $2^{10} = 1024$ iterations of the hash --- enough to make the beacon computationally irreversible.

\begin{highlight}
\textbf{Why a random beacon?} If a contributor could predict the beacon before their contribution, they might be able to bias the parameters. An unpredictable, publicly verifiable source of randomness prevents this. Janus Wallet uses the \texttt{policy\_id} under which all account tokens are minted as its beacon value.
\end{highlight}

\paragraph{7. Prepare Phase 2.}

\begin{verbatim}
snarkjs powersoftau prepare phase2 pot14_beacon.ptau pot14_final.ptau -v
\end{verbatim}

This pre-computes the encrypted Lagrange polynomial evaluations needed in Phase 2.

\paragraph{8. Final verification.}

\begin{verbatim}
snarkjs powersoftau verify pot14_final.ptau
\end{verbatim}

Unlike step 5, this run will \emph{not} warn about missing phase 2 values --- that means the file is ready.

\subsection*{Phase 2: Circuit-Specific Setup}

\paragraph{9--10. Compile the circuit and generate the initial zkey.}

\begin{verbatim}
snarkjs groth16 setup circuit.r1cs pot14_final.ptau circuit_0000.zkey
\end{verbatim}

\begin{remark}
Do not use \texttt{circuit\_0000.zkey} in production. It has zero contributions and is not safe.
\end{remark}

\paragraph{11--12. Add contributions to Phase 2.}

\begin{verbatim}
snarkjs zkey contribute circuit_0000.zkey circuit_0001.zkey \
    --name="1st Contributor Name" -v

snarkjs zkey contribute circuit_0001.zkey circuit_0002.zkey \
    --name="Second contribution" -v -e="Another random entropy"
\end{verbatim}

\paragraph{13. (Optional) Third-party Phase 2 contribution.}

\begin{verbatim}
snarkjs zkey export bellman circuit_0002.zkey challenge_phase2_0003
snarkjs zkey bellman contribute bls12381 challenge_phase2_0003 \
    response_phase2_0003 -e="some random text"
snarkjs zkey import bellman circuit_0002.zkey response_phase2_0003 \
    circuit_0003.zkey -n="Third contribution name"
\end{verbatim}

\paragraph{14. Verify the zkey.}

\begin{verbatim}
snarkjs zkey verify circuit.r1cs pot14_final.ptau circuit_0003.zkey
\end{verbatim}

Output: \texttt{[INFO] snarkJS: ZKey Ok!}

\paragraph{15. Apply the random beacon and export the verification key.}

\begin{verbatim}
snarkjs zkey beacon circuit_0003.zkey circuit_final.zkey \
    0102030405060708090a0b0c0d0e0f101112131415161718191a1b1c1d1e1f \
    10 -n="Final Beacon phase2"

snarkjs zkey verify circuit.r1cs pot14_final.ptau circuit_final.zkey

snarkjs zkey export verificationkey circuit_final.zkey verification_key.json
\end{verbatim}

\texttt{verification\_key.json} contains the public parameters that will be embedded in the Aiken smart contract (as an on-chain token) and used to verify every proof.

% ─────────────────────────────────────────────────────────────────────────────
\section{Generating and Verifying a Proof} \label{sec:proof-gen}

\subsection*{Prepare Inputs}

Create an \texttt{input.json} with the circuit's signals. Note that large integers must be quoted strings to avoid JSON precision loss:

\begin{verbatim}
{
    "userId": "52435875175126190479447740508185965837690552500527637822...",
    "credentialHash": "17469782945530594987388377835453341901430155623074...",
    "challenge": "17469782945530594987388377835453341901430155623074104...",
    "challengeFlag": "1",
    "circuitHash": "24883872159189864133794802823127472657146334491119345...",
    "pwd": "52435875175126190479447740508185965837690552500527637822..."
}
\end{verbatim}

\subsection*{Calculate the Witness}

The witness is the complete set of wire values for all signals in the circuit given your specific inputs.

\begin{verbatim}
snarkjs wtns calculate circuit_js/circuit.wasm input.json witness.wtns
snarkjs wtns check circuit.r1cs witness.wtns
\end{verbatim}

A valid witness prints \texttt{WITNESS IS CORRECT}.

\subsection*{Generate the Proof}

\begin{verbatim}
snarkjs groth16 prove circuit_final.zkey witness.wtns proof.json public.json
\end{verbatim}

This produces:
\begin{itemize}
  \item \texttt{proof.json} --- the Groth16 proof ($\pi_A$, $\pi_B$, $\pi_C$ elliptic curve points).
  \item \texttt{public.json} --- the values of the public inputs (visible to the verifier).
\end{itemize}

Or in a single step:

\begin{verbatim}
snarkjs groth16 fullprove input.json circuit_js/circuit.wasm \
    circuit_final.zkey proof.json public.json
\end{verbatim}

\subsection*{Verify Off-Chain}

\begin{verbatim}
snarkjs groth16 verify verification_key.json public.json proof.json
\end{verbatim}

Output: \texttt{[INFO] snarkJS: OK!}

% ─────────────────────────────────────────────────────────────────────────────
\section{Preparing the Verification Key for Cardano} \label{sec:vk-cardano}

Cardano's BLS12-381 primitives work exclusively with \textbf{compressed} elliptic curve points. The \texttt{verification\_key.json} produced by snarkjs uses \emph{uncompressed} points. Before moving on-chain you must compress all $\mathbb{G}_1$ and $\mathbb{G}_2$ points: \texttt{vk\_alpha\_1}, \texttt{vk\_beta\_2}, \texttt{vk\_gamma\_2}, \texttt{vk\_delta\_2}, and every point in the \texttt{IC} array.

Janus Wallet provides a script for this:

\begin{verbatim}
node --no-warnings --experimental-specifier-resolution=node \
     --loader ts-node/esm compress_zk_verification_key.ts \
     janus-wallet/verification_key.json \
     janus-wallet/compressed_verification_key.json
\end{verbatim}

Run this from \texttt{src/zkproof}. The output \texttt{compressed\_verification\_key.json} is what goes on-chain.

% ─────────────────────────────────────────────────────────────────────────────
\section{On-Chain Verification in Aiken} \label{sec:aiken-verification}

The verification key is minted as a token on Cardano using Janus Wallet's circuit minting endpoint. This makes the parameters immutable and publicly auditable --- anyone can inspect the token and verify the ceremony.

\subsection*{Minting the Circuit Token}

\begin{verbatim}
POST http://localhost:3001/circuits
{
    "token_name": "Circuit#000",
    "version": 0,
    "nonce": "some random value",
    "circuit": {
        "vk_alpha1":  "8f944ea2...",
        "vk_beta2":   "81dfe069...",
        "vk_gamma2":  "93e02b60...",
        "vk_delta2":  "910d1b3d...",
        "vk_ic": [
            "8d043ff5...",
            "86fa74a9...",
            ...
        ]
    }
}
\end{verbatim}

The endpoint mints the token under a \texttt{policy\_id} derived from the nonce. The nonce can be the beacon value used in Phase 2, linking the on-chain token cryptographically to the ceremony transcript.

\subsection*{How the Aiken Validator Works}

The Aiken validator performs a Groth16 pairing check using the BLS12-381 primitives:

\begin{enumerate}
  \item It reads the compressed verification key from the circuit token in the reference input.
  \item It reads the proof ($\pi_A$, $\pi_B$, $\pi_C$) from the redeemer.
  \item It reads the public inputs (userId, credentialHash, challenge, challengeFlag, circuitHash) from the datum.
  \item It computes the pairing equation and calls \texttt{final\_verify} to check the proof.
  \item If the proof is valid, the UTxO is unlocked. If not, the transaction fails.
\end{enumerate}

The Aiken test suite from Janus Wallet shows all passing scenarios:

\begin{verbatim}
aiken check -t silent .

+- tests/verify ----------------------------------------------
| PASS [mem:  89791, cpu: 2541340679] verify_zk_proof
| PASS [mem:  89791, cpu: 2541340679] verify_zk_proof_1
| PASS [mem:  63964, cpu: 1355803465] verify_zk_proof_invalid_user
| PASS [mem:  63964, cpu: 1355812874] verify_zk_proof_invalid_circuit_hash
| PASS [mem:  63964, cpu: 1355812874] verify_zk_proof_invalid_challenge
| PASS [mem:  63964, cpu: 1355812874] verify_zk_proof_invalid_proof
+----------------------------------------- 11 tests | 11 passed | 0 failed
\end{verbatim}

Notice the CPU budget for \texttt{verify\_zk\_proof}: approximately 2.5 billion units. Groth16 verification on BLS12-381 is one of the most expensive operations you can perform on-chain, but it is within Cardano's per-transaction budget --- which is why Plutus V3 added these primitives in the first place.

% ─────────────────────────────────────────────────────────────────────────────
\section{Janus Wallet: Putting It All Together} \label{sec:janus}

Janus Wallet (\texttt{https://github.com/leobel/janus-wallet}) is the reference implementation of a ZK smart contract wallet on Cardano. Instead of signing transactions with a private key, the user proves knowledge of their password via a Groth16 circuit. The private key never exists --- there is nothing to lose or steal.

\subsection*{Architecture}

\begin{itemize}
  \item \textbf{Circuit} --- the Circom Authenticate template described above.
  \item \textbf{Off-chain server} --- a Node.js HTTP API that builds, signs (via ZK proof), and submits transactions.
  \item \textbf{On-chain validator} --- an Aiken smart contract that verifies the Groth16 proof against the circuit token.
  \item \textbf{PostgreSQL} --- stores user accounts, circuit data, and UTxO state.
\end{itemize}

\subsection*{Running the Server}

Create a \texttt{.env} file:

\begin{verbatim}
PORT=3001
DB_HOST=localhost
DB_PORT=5432
DB_USERNAME=janus
DB_PASSWORD=123456
DB_NAME=janus_wallet
CARDANO_NETWORK=Preview
BLOCKFROST_API_URL=https://cardano-preview.blockfrost.io/api/v0
BLOCKFROST_API_KEY=preview...
WALLET_PRV_KEY=ed25519_sk...
COLLATERAL_PRV_KEY=ed25519_sk...
COLLATERAL_ADDRESS=addr_test...
\end{verbatim}

Run database migrations and start the server:

\begin{verbatim}
npm run migrate
npm start
\end{verbatim}

\begin{remark}
The signing logic that produces the ZK proof depends on \texttt{circom} being accessible from the process running the server. Verify with \texttt{circom --version}. If the binary is at \texttt{\$HOME/.cargo/bin} and not on your PATH, copy it to \texttt{/usr/local/bin}.
\end{remark}

\subsection*{User Lifecycle}

\paragraph{Register a user.}

\begin{verbatim}
POST http://localhost:3001/wallets/{user_name}
{
    "hash":     "42b79954...",  // Poseidon(password, username)
    "kdf_hash": "243262..."     // bcrypt(password, salt)
}
\end{verbatim}

The server generates a smart contract address unique to this user and mints a profile NFT to it. The \texttt{hash} is the on-chain credential; the \texttt{kdf\_hash} lets the server authenticate requests without storing the plaintext password.

\paragraph{Send funds (three steps).}

\begin{verbatim}
// 1. Build the unsigned transaction
POST http://localhost:3001/wallets/{user_id}/build
{ "amount": 5000000, "receive_address": "addr_test1..." }
// Response: { "tx": "84a900..." }

// 2. Sign with the ZK proof
POST http://localhost:3001/wallets/{user_id}/sign
{ "pwd": "supersecret", "tx": "84a900..." }

// 3. Submit
POST http://localhost:3001/wallets/{user_id}/send
{
    "redeemers": ["d8799f..."],  // ZK proof as CBOR
    "tx": "84a900..."
}
\end{verbatim}

The \texttt{redeemers} array contains the Groth16 proof serialised as CBOR. Additional redeemers in the array are bypass markers for UTxOs from the same wallet that do not require a fresh proof.

\paragraph{Stake, delegate, and withdraw.}

Janus Wallet supports the full Cardano staking lifecycle through the same three-step pattern (build $\to$ sign $\to$ send):

\begin{verbatim}
POST /wallets/{user_id}/pools/{pool_id}/registerAndDelegate
POST /wallets/{user_id}/pools/{pool_id}/delegate
POST /wallets/{user_id}/dreps/{drep_id}/delegate
POST /wallets/{user_id}/withdraw  { "amount": 1234567 }
\end{verbatim}

\begin{highlight}
\textbf{The key insight:} from the blockchain's perspective, every Janus transaction is a normal smart contract spend. The only thing that makes it a ZK wallet is the redeemer --- a Groth16 proof that the user knows their password, checked on-chain by the Aiken validator in a single pairing verification.
\end{highlight}

\subsection*{Full Test Results}

Janus Wallet ships with 29 Aiken tests covering minting, spending, challenge construction, and ZK verification --- all passing with zero errors:

\begin{verbatim}
Summary 29 checks, 0 errors, 19 warnings
\end{verbatim}

The warnings are informational (unused variables in test helpers) and do not affect correctness.

This is the complete picture of ZK on Cardano today: Circom defines the computation, snarkjs runs the ceremony and generates proofs, and Aiken verifies them on-chain using native BLS12-381 pairing operations. In the next chapters we will build our own circuit and validator from scratch.
